\documentclass[11pt]{article}
\usepackage[margin=1in]{geometry}
\usepackage{amsmath,amssymb,amsthm,mathtools}
\usepackage{graphicx}
\usepackage[hidelinks]{hyperref}
\usepackage{microtype}

\newtheorem{theorem}{Theorem}
\newtheorem{lemma}{Lemma}
\newcommand{\R}{\mathbb R}
\newcommand{\eps}{\varepsilon}
\newcommand{\Lp}{L^p[0,1]}
\newcommand{\Xq}{X_q}

\title{A logarithmic endpoint partial result for the\\
continuous Bennett--Carl conjecture}
\author{Run \texttt{fa\_banach\_001}, lane 13}
\date{21 July 2026}

\begin{document}
\maketitle

\paragraph{Status.}
Candidate partial result, likely valid.  The $(p,1)$-summing conjecture itself
remains open.

\section{Source problem}

Astashkin, Le\'snik and Wojciechowski~\cite{ALW} define, for $1<p<2$,
\[
  X_p=\Lambda_{p,W_p},\qquad
  W_p(t)=\log^{,1-p}(e/t),
\]
and prove that the inclusion $J:X_p\to L^p[0,1]$ is $(q,1)$-absolutely
summing for every $p<q<2$.  Their Conjecture~3.4 asks whether $J$ is
$(p,1)$-absolutely summing.

\begin{center}
\includegraphics[width=0.98\textwidth]{figures/open_problem_crop.png}
\end{center}

The crop is from page 12 of the source PDF.  In finite-family form the
conjecture asks for a constant $C_p$ such that
\begin{equation}\label{eq:conj}
 \left(\sum_i\|f_i\|_p^p\right)^{1/p}
 \le C_p\max_{\eps_i=\pm1}
       \left\|\sum_i\eps_i f_i\right\|_{X_p}.
\end{equation}

For a Young function $\Phi$, write $\ell_\Phi$ for the Luxemburg sequence
space.  We say that $T:E\to F$ is $(\Phi,1)$-summing if
\[
 \big\|(\|Tx_i\|_F)_i\big\|_{\ell_\Phi}
 \le C\sup_{x^*\in B_{E^*}}\sum_i|x^*(x_i)|
\]
for all finite families.  The quantity on the right equals
$\max_{\eps_i=\pm1}\|\sum_i\eps_i x_i\|_E$.

\section{Proof intuition}

The source proof loses the endpoint only in the embedding of the exponential
Orlicz space $\operatorname{Exp}L^{p'}$ into $X_q$.  Its norm can be computed
from rearrangements: the $q$-th power grows like $(q-p)^{-1}$.  All other
constants can be kept uniform while $q\downarrow p$.  Consequently
\[
  \pi_{q,1}(J)^q\lesssim_p (q-p)^{-1}.
\]
Integrating this family of estimates with weight
$(q-p)^{\beta-1}$ turns $u^{p+q-p}$ into
$u^p/\log^\beta(e/u)$.  This produces a single logarithmic endpoint space
for every $\beta>1$.  The integral diverges at $\beta=1$, so the mechanism
does not prove the original conjecture.

\section{The partial theorem}

For $\beta>1$, let $\Phi_\beta$ be any Young function such that
\begin{equation}\label{eq:phi}
 \Phi_\beta(u)\asymp_{p,\beta}
 \frac{u^p}{\log^\beta(e/u)}\qquad(0<u\le u_0).
\end{equation}
Such a Young function is obtained by convexifying the displayed function near
zero and continuing it convexly; equivalent choices give the same sequence
space up to equivalent norms.

\begin{theorem}\label{thm:main}
Let $1<p<2$ and $\beta>1$.  The inclusion
$J:X_p\to L^p[0,1]$ is $(\Phi_\beta,1)$-absolutely summing.  Thus there is
$C_{p,\beta}<\infty$ such that every finite family $(f_i)\subset X_p$ obeys
\begin{equation}\label{eq:orlicz}
 \big\|(\|f_i\|_p)_i\big\|_{\ell_{\Phi_\beta}}
 \le C_{p,\beta}\max_{\eps_i=\pm1}
 \left\|\sum_i\eps_i f_i\right\|_{X_p}.
\end{equation}
Moreover,
\[
 \ell^p\subsetneq\ell_{\Phi_\beta}\subsetneq\ell^q
 \quad\text{for every }q>p.
\]
In particular, \eqref{eq:orlicz} is a genuine endpoint-scale strengthening of
the separate $(q,1)$ estimates in~\cite{ALW}, but it is weaker than
\eqref{eq:conj}.
\end{theorem}

\section{Proof}

We first quantify the source argument.  Put $p'=p/(p-1)$ and
$E=\operatorname{Exp}L^{p'}$.

\begin{lemma}[Uniform Rademacher lower estimate]\label{lem:rad}
There exist $\eta>0$ and $c_p>0$ such that, for every
$q\in(p,p+\eta)$ and every finite scalar family $(a_i)$,
\[
 \left\|\sum_i a_i r_i\right\|_{X_q}
 \ge c_p\left(\sum_i|a_i|^q\right)^{1/q}.
\]
\end{lemma}

\begin{proof}
Reorder so that $|a_1|\ge\cdots\ge|a_n|$ and set
$A_k=\sum_{i\le k}|a_i|$.  The first $k$ signs align with the coefficients
with probability $2^{-k}$.  Conditional on this event, the remaining signed
sum is symmetric and is nonnegative with probability at least $1/2$.
Therefore
\[
 \left(\sum_i a_i r_i\right)^*(2^{-k-1})\ge A_k.
\]
Writing $W_q(t)=\log^{1-q}(e/t)$ and integrating on the adjacent dyadic
intervals gives
\[
 \left\|\sum_i a_i r_i\right\|_{X_q}^q
 \ge \sum_{k=1}^n A_k^q
 \bigl[W_q(2^{-k-1})-W_q(2^{-k-2})\bigr].
\]
By the mean value theorem, uniformly for $q$ in a sufficiently small compact
interval above $p$, the bracket is at least $c_p^q(k+2)^{-q}$.  Since
$A_k/k\ge |a_k|$, the last display is at least
$c_p^q\sum_k|a_k|^q$, after changing $c_p$ by a fixed factor.
\end{proof}

\begin{lemma}[Quantitative near-endpoint estimate]\label{lem:q}
After possibly decreasing $\eta$, for $p<q<p+\eta$,
\begin{equation}\label{eq:qbound}
 \pi_{q,1}(J)\le C_p(q-p)^{-1/q}.
\end{equation}
\end{lemma}

\begin{proof}
The elementary rearrangement estimate
\[
 \|f\|_{X_p}^p\ge f^*(t)^pW_p(t)
\]
implies
$f^*(t)\le\|f\|_{X_p}\log^{1/p'}(e/t)$, and hence
$X_p\hookrightarrow E$ with a constant depending only on $p$.
Conversely, the standard rearrangement bound in $E$ gives
$f^*(t)\le C_p\|f\|_E\log^{1/p'}(e/t)$.  Since
\[
 dW_q(t)=\frac{q-1}{t\log^q(e/t)}\,dt,
\]
we obtain
\begin{align*}
 \|f\|_{X_q}^q
 &\le C_p^q\|f\|_E^q(q-1)
   \int_0^1\frac{\log^{q/p'}(e/t)}
                    {t\log^q(e/t)}\,dt\\
 &=C_p^q\|f\|_E^q(q-1)
   \int_1^\infty u^{-q/p}\,du\\
 &=C_p^q\|f\|_E^q\frac{p(q-1)}{q-p}.
\end{align*}
Thus
\begin{equation}\label{eq:embed}
 \|E\hookrightarrow X_q\|
 \le C_p(q-p)^{-1/q}
\end{equation}
uniformly in the stated interval.

Let $(f_i)\subset X_p$.  The $q$-concavity inequality in $L^p$ and
Lemma~\ref{lem:rad}, applied pointwise, yield
\[
 \left(\sum_i\|f_i\|_p^q\right)^{1/q}
 \le C_p\left\|\left\|\sum_i r_i(s)f_i(t)
                  \right\|_{X_q(s)}\right\|_{L^p(t)}.
\]
Use \eqref{eq:embed}, followed by the transposition estimate
\[
 \mathcal T:L^\infty(E)\longrightarrow L^p(E)
\]
from Lemma~3.1 of~\cite{ALW}, and finally $X_p\hookrightarrow E$.  We get
\begin{align*}
 \left(\sum_i\|f_i\|_p^q\right)^{1/q}
 &\le C_p(q-p)^{-1/q}
 \left\|\left\|\sum_i r_i(t)f_i(s)
                  \right\|_{X_p(s)}\right\|_{L^\infty(t)}\\
 &=C_p(q-p)^{-1/q}
 \max_{\eps_i=\pm1}\left\|\sum_i\eps_i f_i\right\|_{X_p}.
\end{align*}
This is \eqref{eq:qbound}.
\end{proof}

\begin{proof}[Proof of Theorem~\ref{thm:main}]
Normalize
\[
 A:=\max_{\eps_i=\pm1}\left\|\sum_i\eps_i f_i\right\|_{X_p}=1,
 \qquad y_i=\|f_i\|_p.
\]
Lemma~\ref{lem:q} says, for $0<s<\eta$,
\begin{equation}\label{eq:moments}
 \sum_i y_i^{p+s}\le \frac{C_p^{p+s}}{s}.
\end{equation}
Also $\sup_i y_i\le C_p$: prescribed signs on one coordinate can be extended
to all coordinates without decreasing the norm, and $X_p\hookrightarrow L^p$.

Choose $K\ge2C_p$ and put $z_i=y_i/K\le1/2$.  Multiply
\eqref{eq:moments} by $K^{-(p+s)}s^{\beta-1}$ and integrate over
$0<s<\eta$.  Enlarging $K$ by a fixed $p$-dependent factor if necessary,
\begin{equation}\label{eq:upperint}
 \sum_i\int_0^\eta z_i^{p+s}s^{\beta-1}\,ds
 \le C_{p,\beta}K^{-p}\int_0^\eta s^{\beta-2}\,ds
 <\infty,
\end{equation}
where the last integral is finite exactly because $\beta>1$.

For $0<z\le1/2$, put $L=\log(1/z)$.  Integrating only over
$0<s<\min(\eta,L^{-1})$ shows
\[
 \int_0^\eta z^{p+s}s^{\beta-1}\,ds
 =z^p\int_0^\eta e^{-sL}s^{\beta-1}\,ds
 \ge c_{p,\beta,\eta}\frac{z^p}{\log^\beta(e/z)}.
\]
Together with \eqref{eq:upperint}, this proves
\[
 \sum_i\Phi_\beta(y_i/K)\le C'_{p,\beta}K^{-p}.
\]
Increasing $K=K(p,\beta)$ once more makes the right side at most one, which
is precisely the Luxemburg estimate \eqref{eq:orlicz}.  Homogeneity restores
an arbitrary value of $A$.

Finally, as $u\downarrow0$,
$\Phi_\beta(u)=o(u^p)$, whereas $u^q=o(\Phi_\beta(u))$ for every $q>p$.
The standard characteristic-sequence tests give the asserted strict
inclusions.
\end{proof}

\section{Verification, limitations, and novelty status}

\begin{itemize}
\item The proof isolates every dependence on $q-p$.  Lemma~\ref{lem:rad}
is uniform by an explicit dyadic calculation; the sole singular factor is
the integral in \eqref{eq:embed}.
\item The extrapolation argument is written out rather than invoked as a
black box.  It is the same Yano-type mechanism developed in Section~8 of
Kazaniecki--Pawlewicz--Wojciechowski~\cite{KPW}.
\item The argument stops at $\beta>1$.  At $\beta=1$, the upper integral in
\eqref{eq:upperint} is logarithmically divergent.  No assertion is made for
$\beta=1$, and the desired $\Phi(u)=u^p$ endpoint remains open.
\item Exact finite-atomic searches from the preceding attempt found no ratio
above one, including perturbations of the disjoint extremizers.  These tests
support, but do not prove, Conjecture~3.4 and are not used above.
\item A bounded web/arXiv and local-corpus search on 21 July 2026 used the
source title, arXiv id, author combinations, ``continuous Bennett--Carl'',
``endpoint summing'', and logarithmic Lorentz/Orlicz variants.  It found no
later settlement and no explicit statement of Theorem~\ref{thm:main}.
\end{itemize}

\paragraph{Human-review recommendation.}
Promote as a likely-valid partial theorem, not as a solution of Conjecture~3.4.
The most important checks are the uniformity in Lemma~\ref{lem:rad}, the
orientation of the two mixed norms in the transposition step, and the
Luxemburg normalization in the final integral.

\begin{thebibliography}{9}
\bibitem{ALW}
S.~V. Astashkin, K.~Le\'snik, and M.~Wojciechowski,
\emph{Looking for a continuous version of Bennett--Carl theorem},
arXiv:2502.07041 (2025).

\bibitem{KPW}
K.~Kazaniecki, A.~Pawlewicz, and M.~Wojciechowski,
\emph{Factoring the Sobolev embedding operator},
arXiv:2503.20478 (2025), Section~8.
\end{thebibliography}

\end{document}
